\section{Extensions}
\label{section-extensions}

\begin{definition}
\label{definition-extension}
Let $\mathcal{A}$ be an abelian category.
Let $A, B \in \Ob(\mathcal{A})$.
An {\it extension $E$ of $B$ by $A$} is a short
exact sequence
$$
0 \to A \to E \to B \to 0.
$$
A {\it morphism of extensions} between two
extensions $0 \to A \to E \to B \to 0$ and
$0 \to A \to F \to B \to 0$ means a morphism
$f : E \to F$ in $\mathcal{A}$ making the diagram
$$
\xymatrix{
0 \ar[r] &
A \ar[r] \ar[d]^{\text{id}} &
E \ar[r] \ar[d]^f &
B \ar[r] \ar[d]^{\text{id}} &
0 \\
0 \ar[r] &
A \ar[r] &
F \ar[r] &
B \ar[r] &
0
}
$$
commutative.
Thus, the extensions of $B$ by $A$ form a category.
\end{definition}

\noindent
By abuse of language we often omit mention of the
morphisms $A \to E$ and $E \to B$, although they are
definitively part of the structure of an extension.

\begin{definition}
\label{definition-ext-group}
Let $\mathcal{A}$ be an abelian category.
Let $A, B \in \Ob(\mathcal{A})$.
The set of isomorphism classes of extensions
of $B$ by $A$ is denoted
$$
\Ext_\mathcal{A}(B, A).
$$
This is called the {\it $\Ext$-group}.
\end{definition}

\noindent
This definition works, because by our conventions
$\Ob(\mathcal{A})$ is a set, and hence
$\Ext_\mathcal{A}(B, A)$ is a set.
In any of the cases of ``big'' abelian categories
listed in Categories, Remark \ref{categories-remark-big-categories}
one can check by hand that $\Ext_\mathcal{A}(B, A)$
is a set as well. Also, we will see later that this is
always the case when $\mathcal{A}$ has either enough projectives
or enough injectives. Insert future reference here.

\medskip\noindent
Actually we can turn $\Ext_\mathcal{A}(-, -)$ into a
functor
$$
\mathcal{A} \times \mathcal{A}^{opp} \longrightarrow \textit{Sets}, \quad
(A, B) \longmapsto \Ext_\mathcal{A}(B, A)
$$
as follows:
\begin{enumerate}
\item Given a morphism $B' \to B$ and an extension
$E$ of $B$ by $A$ we define $E' = E \times_B B'$.
By Lemma \ref{lemma-cartesian-kernel} we have the following
commutative diagram of short exact sequences
$$
\xymatrix{
0 \ar[r] & A \ar[r] \ar[d] & E' \ar[r] \ar[d] & B' \ar[r] \ar[d] & 0 \\
0 \ar[r] & A \ar[r] & E \ar[r] & B \ar[r] & 0
}
$$
The extension $E'$ is called the {\it pullback of $E$ via
$B' \to B$}.
\item Given a morphism $A \to A'$ and an extension
$E$ of $B$ by $A$ we define $E' = A' \amalg_A E$.
By Lemma \ref{lemma-cartesian-cocartesian}
we have the following commutative diagram
of short exact sequences
$$
\xymatrix{
0 \ar[r] & A \ar[r] \ar[d] & E \ar[r] \ar[d] & B \ar[r] \ar[d] & 0 \\
0 \ar[r] & A' \ar[r] & E' \ar[r] & B \ar[r] & 0
}
$$
The extension $E'$ is called the {\it pushout of $E$ via
$A \to A'$}.
\end{enumerate}
To see that this defines a functor as indicated above
there are several things to verify. First of all
functoriality in the variable $B$ requires that
$(E \times_B B') \times_{B'} B'' = E \times_B B''$
which is a general property of fibre products.
Dually one deals with functoriality in the
variable $A$. Finally, given $A \to A'$ and
$B' \to B$ we have to show that
$$
A' \amalg_A (E \times_B B')
\cong
(A' \amalg_A E)\times_B B'
$$
as extensions of $B'$ by $A'$. Recall that $A' \amalg_A E$
is a quotient of $A' \oplus E$.
Thus the right hand side is a quotient of
$A' \oplus E \times_B B'$, and it is straightforward to see that
the kernel is exactly what you need in order to
get the left hand side.

\medskip\noindent
Note that if $E_1$ and $E_2$ are extensions of
$B$ by $A$, then $E_1\oplus E_2$ is an extension
of $B \oplus B$ by $A\oplus A$. We
push out by the sum map $A \oplus A \to A$ and we
pull back by the diagonal map $B \to B \oplus B$ to get
an extension $E_1 + E_2$ of $B$ by $A$.
$$
\xymatrix{
0 \ar[r] &
A \oplus A \ar[r] \ar[d]_\Sigma &
E_1 \oplus E_2 \ar[r] \ar[d] &
B \oplus B \ar[r] \ar[d] &
0 \\
0 \ar[r] &
A \ar[r] &
E' \ar[r] &
B \oplus B \ar[r] &
0\\
0 \ar[r] &
A \ar[r] \ar[u] &
E_1 + E_2 \ar[r] \ar[u] &
B \ar[r] \ar[u]^\Delta &
0
}
$$
The extension $E_1 + E_2$ is called the {\it Baer sum} of the
given extensions.

\begin{lemma}
\label{lemma-baer-sum}
The construction $(E_1, E_2) \mapsto E_1 + E_2$
above defines a commutative group
law on $\Ext_\mathcal{A}(B, A)$ which is
functorial in both variables.
\end{lemma}

\begin{proof}
Omitted.
\end{proof}

\begin{lemma}
\label{lemma-six-term-sequence-ext}
Let $\mathcal{A}$ be an abelian category.
Let $0 \to M_1 \to M_2 \to M_3 \to 0$ be a short exact sequence
in $\mathcal{A}$.
\begin{enumerate}
\item There is a canonical six term exact sequence of abelian groups
$$
\xymatrix{
0 \ar[r] &
\Hom_\mathcal{A}(M_3, N) \ar[r] &
\Hom_\mathcal{A}(M_2, N) \ar[r] &
\Hom_\mathcal{A}(M_1, N) \ar[lld] \\
& \Ext_\mathcal{A}(M_3, N) \ar[r] &
\Ext_\mathcal{A}(M_2, N) \ar[r] &
\Ext_\mathcal{A}(M_1, N)
}
$$
for all objects $N$ of $\mathcal{A}$, and
\item there is a canonical six term exact sequence of abelian groups
$$
\xymatrix{
0 \ar[r] &
\Hom_\mathcal{A}(N, M_1) \ar[r] &
\Hom_\mathcal{A}(N, M_2) \ar[r] &
\Hom_\mathcal{A}(N, M_3) \ar[lld] \\
& \Ext_\mathcal{A}(N, M_1) \ar[r] &
\Ext_\mathcal{A}(N, M_2) \ar[r] &
\Ext_\mathcal{A}(N, M_3)
}
$$
for all objects $N$ of $\mathcal{A}$.
\end{enumerate}
\end{lemma}

\begin{proof}
Omitted. Hint: The boundary maps are defined using either the pushout
or pullback of the given short exact sequence.
\end{proof}






